% 星际行星（综述）
% license CCBYSA3
% type Wiki

本文根据 CC-BY-SA 协议转载翻译自维基百科\href{https://en.wikipedia.org/wiki/Rogue_planet}{相关文章}。

流浪行星（rogue planet），亦称自由漂浮行星（free-floating planet，FFP）或孤立的行星质量天体（isolated planetary-mass object，iPMO），是指一种具有行星质量、但在引力上不与任何恒星或棕矮星发生束缚关系的星际天体$^{[1][2][3][4]}$。

流浪行星可能起源于其形成所在的行星系统，并在随后被抛射出去；它们也可能在行星系统之外独立形成。仅银河系中就可能存在数十亿到数万亿颗流浪行星，这一数量范围有望由即将投入使用的南希·格雷斯·罗曼空间望远镜进一步加以精确限定$^{[5][6]}$。流浪行星进入太阳系的概率极低，更不用说对地球生命构成直接威胁——在未来一千年内发生这种情况的概率约为一万亿分之一$^{[7]}$。

一些行星质量天体可能以类似恒星的方式形成，国际天文学联合会已提议将此类天体称为亚棕矮星（sub-brown dwarfs）$^{[8]}$。一个可能的例子是 Cha~110913$-$773444，它可能是在被抛射后成为一颗流浪行星，也可能是独立形成并成为一颗亚棕矮星$^{[9]}$。
\subsection{术语}
最早的两篇发现论文分别使用了“孤立的行星质量天体”（iPMO）$^{[10]}$和“自由漂浮行星”（FFP）$^{[11]}$这两个名称。大多数天文学论文使用其中之一$^{[12][13][14]}$。“流浪行星”这一术语更常用于引力微透镜研究，而这类研究也经常使用 FFP 这一术语$^{[15][16]}$。面向公众的新闻稿则可能采用不同的名称。例如，2021 年发现至少 70 颗 FFP 时，不同新闻稿分别使用了“流浪行星”$^{[17]}$、“无恒星行星”$^{[18]}$、“游荡行星”$^{[19]}$以及“自由漂浮行星”$^{[20]}$等称呼。
\subsection{发现}
孤立的行星质量天体（isolated planetary-mass objects，iPMO）最早于 2000 年由英国研究团队 Lucas \& Roche 使用英国红外望远镜（UKIRT）在猎户座星云中发现$^{[11]}$。同年，西班牙研究团队 Zapatero Osorio 等人利用凯克望远镜的光谱观测，在 $\sigma$~猎户座星团中发现了 iPMO$^{[10]}$。对猎户座星云中这些天体的光谱研究于 2001 年发表$^{[21]}$。目前，这两个欧洲研究团队均因其准同步发现而受到认可$^{[22]}$。1999 年，日本研究团队 Oasa 等人在变色龙座 I 云中发现了一些天体$^{[23]}$，这些天体随后在 2004 年由美国研究团队 Luhman 等人通过光谱观测得到了确认$^{[24]}$。
\subsection{观测}
\begin{figure}[ht]
\centering
\includegraphics[width=6cm]{./figures/35a477b3258b7d43.png}
\caption{2021 年在上天蝎座与蛇夫座之间的区域发现了 115 个潜在的流浪行星} \label{fig_RoPl_1}
\end{figure}
发现自由漂浮行星主要有两种技术：直接成像和引力微透镜。
\subsubsection{引力微透镜}
日本大阪大学的天体物理学家角广隆宏（Takahiro Sumi）及其同事，隶属于天体物理学引力微透镜观测计划（Microlensing Observations in Astrophysics）和光学引力透镜实验（Optical Gravitational Lensing Experiment）合作组，于 2011 年发表了他们关于引力微透镜的研究成果。他们使用位于新西兰约翰山天文台的 1.8 米（5 英尺 11 英寸）MOA-II 望远镜，以及位于智利拉斯坎帕纳斯天文台的华沙大学 1.3 米（4 英尺 3 英寸）望远镜，对银河系中约 5000 万颗恒星进行了观测。他们共发现了 474 次微透镜事件，其中有 10 次事件的持续时间足够短，表明它们可能是质量约为木星质量、且在其附近没有伴随恒星的行星。研究人员根据这些观测结果估计，在银河系中，每一颗恒星大约对应近两颗木星质量的流浪行星$^{[25][26][27]}$。另有一项研究提出了一个大得多的数量估计，认为银河系中的流浪行星数量可能是恒星数量的多达 100{,}000 倍，不过该研究所涵盖的假想天体质量远小于木星$^{[28]}$。2017 年，华沙大学天文台的 Przemek~Mróz 及其同事在统计样本数量为 2011 年研究六倍的基础上开展的一项研究表明，银河系中每一颗主序星所对应的木星质量自由漂浮行星或宽轨道行星的上限为 0.25 颗$^{[29]}$。

2020 年 9 月，天文学家利用引力微透镜技术首次报告探测到一颗地球质量的流浪行星（命名为 OGLE-2016-BLG-1928），该天体不与任何恒星发生引力束缚，并在银河系中自由漂浮$^{[16][30][31]}$。
\subsubsection{直接成像}
\begin{figure}[ht]
\centering
\includegraphics[width=6cm]{./figures/250326c529d9a895.png}
\caption{使用斯皮策空间望远镜观测到的低温行星质量天体 WISE~J0830+2837（以橙色天体标示）。其温度约为 300--350~K（27--77\,$^\circ$C；80--170\,$^\circ$F）。} \label{fig_RoPl_2}
\end{figure}
微透镜行星只能通过微透镜事件本身进行研究，这使得对行星性质的刻画变得困难。因此，天文学家转而研究通过直接成像方法发现的孤立行星质量天体（iPMO）。例如，要确定一颗棕矮星或 iPMO 的质量，需要其光度和年龄等信息$^{[32]}$。然而，确定低质量天体的年龄已被证明是困难的。因此，绝大多数 iPMO 被发现于附近年轻的恒星形成区中，这些区域的年龄为天文学家所知。这些天体的年龄小于 200~Myr，质量较大（$>5\,M_J$）$^{[4]}$，并且属于 $L$ 矮星和 $T$ 矮星$^{[33][34]}$。然而，也存在一个规模虽小但正在增长的低温且年老的 $Y$ 矮星样本，其估计质量为 $8$--$20\,M_J$ $^{[35]}$。光谱型为 $Y$ 的附近流浪行星候选体包括距离为 $7.27\pm0.13$ 光年的 WISE~0855$-$0714$^{[36]}$。如果能够通过更精确的测量对这一 Y 矮星样本进行刻画，或找到更好地确定其年龄的方法，那么年老且低温的 iPMO 数量很可能会显著增加。

最早的 iPMO 于 21 世纪初通过对年轻恒星形成区的直接成像而被发现$^{[37][10][21]}$。这些通过直接成像发现的 iPMO 很可能以类似恒星的方式形成（有时被称为亚棕矮星）。也可能存在以行星方式形成、随后被抛射出去的 iPMO。然而，这类天体在运动学上应当不同于其诞生的恒星形成区，不应被环星盘所包围，并且应具有较高的金属丰度$^{[22]}$。在年轻恒星形成区中发现的 iPMO 中，没有任何一个相对于其恒星形成区表现出较高的速度。对于年老的 iPMO，低温天体 WISE~J0830+2837$^{[38]}$ 的切向速度 $V_{\mathrm{tan}}$ 约为 $100\,\mathrm{km/s}$，这一数值虽高，但仍与其在银河系中形成的情形相一致。对于 WISE~1534$-$1043$^{[39]}$，由于其约为 $200\,\mathrm{km/s}$ 的高 $V_{\mathrm{tan}}$，一种替代情景将其解释为被抛射出的系外行星，但其颜色却表明它是一颗年老、低金属丰度的棕矮星。大多数研究大质量 iPMO 的天文学家认为，它们代表了恒星形成过程的低质量端$^{[22]}$。

天文学家利用赫歇尔空间天文台和甚大望远镜观测了一颗非常年轻的自由漂浮行星质量天体 OTS~44，并证明了典型的类恒星形成模式所具有的过程同样适用于质量低至数个木星质量的孤立天体。赫歇尔的远红外观测显示，OTS~44 被一个质量至少为 10 个地球质量的盘所包围，因此最终可能形成一个大型卫星系统$^{[40]}$。利用甚大望远镜上的 SINFONI 光谱仪对 OTS~44 进行的光谱观测揭示，该盘正在积极地吸积物质，这与年轻恒星周围的盘相似$^{[40]}$。
\subsubsection{双星}
\begin{figure}[ht]
\centering
\includegraphics[width=6cm]{./figures/2fb0b35b79287356.png}
\caption{2MASS~J1119–1137AB，即首个被发现的行星质量双星系统，位于 TW~Hydrae 协会中} \label{fig_RoPl_3}
\end{figure}
\begin{figure}[ht]
\centering
\includegraphics[width=6cm]{./figures/4a9c799a9b362b83.png}
\caption{JuMBO~29，一个候选的 $12.5+3\,M_J$ 行星质量双星系统，分离距离为 135~AU，位于猎户座星云中} \label{fig_RoPl_4}
\end{figure}
首个被分辨开的行星质量双星的发现对象是 2MASS~J1119–1137AB。然而，目前还已知其他一些双星系统，例如 2MASS~J1553022+153236AB$^{[41][42]}$、WISE~1828+2650、WISE~0146+4234、WISE~J0336–0143（也可能是一个由棕矮星与行星质量天体组成的双星系统（BD+PMO））、NIRISS-NGC1333-12$^{[43]}$，以及 Zhang 等人发现的若干天体$^{[42]}$。

在猎户座星云中，人们发现了由 40 个宽双星系统和 2 个三体系统构成的一群天体。这一发现之所以令人惊讶，有两个原因：其一，针对棕矮星双星的理论趋势预测，随着质量降低，低质量天体之间的分离距离应当减小；其二，还预测双星比例会随质量降低而减少。这些双星被命名为木星质量双星天体（Jupiter-mass Binary Objects，JuMBOs）；它们至少占所有 iPMO 的 9\%，其分离距离小于 340~AU$^{[44]}$。目前尚不清楚这些 JuMBO 的形成机制，但一项深入研究认为它们是在原地形成的，方式类似于恒星$^{[45]}$。如果它们确实以恒星的方式形成，那么必然存在某种尚未明确的“额外因素”使其得以形成；如果它们以行星的方式形成并随后被抛射出去，则需要解释为何这些双星在抛射过程中并未解体。未来利用 JWST 的观测或许能够判定这些天体究竟是以被抛射行星的方式形成，还是以恒星的方式形成$^{[44]}$。Kevin~Luhman 对 NIRCam 数据进行了重新分析，发现大多数 JuMBO 并未出现在他所选取的亚恒星天体样本中。此外，其颜色与被红化的背景源或低信噪比源一致。他仅认为 JuMBO~29 是一个可靠的行星质量双星系统候选体$^{[46]}$。
\subsubsection{已知 iPMO 的总数}
目前已知的候选 iPMO 很可能有数百个$^{[47][44]}$，其中有一百多个天体已获得光谱观测结果$^{[48][49][50]}$，此外还有数量虽少但正在增加的通过引力微透镜发现的候选体。一些大型巡天项目包括：

截至 2021 年 12 月，迄今为止规模最大的流浪行星样本被发现，其数量至少为 70 个，在假定不同年龄的情况下最多可达 170 个。这些天体位于上天蝎座与蛇夫座之间的 OB 协会中，质量介于$4$ 至 $13\,M_J$之间，年龄约为 300 万至 1000 万年，最有可能通过气体云的引力坍缩形成，或在原行星盘中形成后因动力学不稳定而被抛射出去$^{[47][17][51][19]}$。随后利用昴星团望远镜和加那利大型望远镜开展的光谱跟进观测表明，该样本的污染率相当低（$\leq6\%$）。其中 16 个年轻天体的质量介于$3$ 至 $14\,M_J$ 之间，确认它们确实属于行星质量天体$^{[50]}$。

2023 年 10 月，利用 JWST 在梯形星团及猎户座星云内部区域又发现了一个更大的行星质量天体候选体样本，总数达到 540 个。这些天体的质量范围为 $13$ 至 $0.6\,M_J$。其中相当一部分天体形成了宽双星系统，这一现象此前并未被预测到$^{[44]}$。
\subsection{形成}
总体而言，孤立行星质量天体（isolated planetary-mass object，iPMO）的形成主要有两种情景：其一，它像行星一样在恒星周围形成，随后被抛射出去；其二，它像低质量恒星或棕矮星一样在孤立环境中形成。这两种形成方式会影响其成分和运动特性$^{[22]}$。

近期研究表明，流浪行星既可能通过恒星形成区内的直接引力坍缩形成，也可能在其诞生的行星系统中形成后被抛射出来，并在随后与已建立的行星系统发生相互作用，从而影响这些系统的轨道结构和整体统计特征。许多此类天体很可能最初起源于行星系统之中，随后在动力学过程中被驱逐，而另一些则可能是在孤立环境中形成的。除了在近距离遭遇或潜在的俘获事件中改变系统稳定性之外，流浪行星还可能输送挥发物，从而增强前生物化学过程，并创造有利于提高生物多样性的条件。这些形成机制、动力学过程、生化作用以及生态效应的综合影响，在塑造系外行星系统的分布与演化方面发挥着重要作用$^{[52]}$。
\subsubsection{类似恒星的形成}
2001 年的模型认为，质量至少为一个木星质量的天体可以通过分子云的坍缩与碎裂过程形成$^{[53]}$。JWST 之前的观测表明，质量低于$3$--$5\,M_J$ 的天体不太可能独立形成$^{[4]}$。2023 年，利用 JWST 对梯形星团的观测显示，质量低至 $0.6\,M_J$ 的天体也可能独立形成，并不需要一个陡峭的质量截断$^{[44]}$。一种称为小球状体（globulettes）的特殊云核被认为是棕矮星和行星质量天体的诞生地。小球状体可见于玫瑰星云和 IC~1805 中$^{[54]}$。有时，年轻的 iPMO 仍被一个可能形成系外卫星的盘所包围。由于这种系外卫星围绕其宿主行星的轨道非常紧密，其发生凌星现象的概率高达$10$--$15\%$，$^{[55]}$。

\textbf{盘}

一些非常年轻的恒星形成区（通常年龄小于 500 万年）有时包含具有红外过量并显示吸积迹象的孤立行星质量天体。其中最著名的是位于变色龙座 I 的 iPMO——OTS~44，其被发现具有一个盘。变色龙座 I 和 II 中还存在其他带盘的 iPMO 候选体$^{[56][57][33]}$。其他包含带盘或吸积 iPMO 的恒星形成区包括：天狼座 I$^{[57]}$、蛇夫座 $\rho$ 云复合体$^{[58]}$、$\sigma$~猎户座星团$^{[59]}$、猎户座星云$^{[60]}$、金牛座$^{[58][61]}$、NGC~1333$^{[62]}$ 以及 IC~348$^{[63]}$。利用 ALMA 对棕矮星和 iPMO 周围盘的大型巡天发现，这些盘的质量不足以形成地球质量的行星，但仍存在这些盘中已形成行星的可能性$^{[58]}$。对红矮星的研究表明，其中一些在相对较老的年龄阶段仍保留富含气体的盘，这类盘被称为“彼得·潘盘”（Peter Pan Disks），而这一趋势可能延伸至行星质量范围。一个彼得·潘盘的例子是年龄约为 45~Myr、质量约为$13.7\,M_J$ 的棕矮星 2MASS~J02265658-5327032，其质量已接近行星质量范围$^{[64]}$。近期对附近行星质量天体 2MASS~J11151597+1937266 的研究发现，该 iPMO 周围存在一个盘，并显示出来自盘的吸积迹象以及红外过量$^{[65]}$。2025 年 5 月，研究人员利用 JWST 发现 Cha~1107-7626 周围的盘中含有烃类物质。Cha~1107-7626（$6--10\,M_J$）是已知具有尘埃盘的最低质量天体之一$^{[66]}$。进一步的 JWST 光谱观测表明，硅酸盐和烃类是行星质量天体盘中的常见成分。这些盘显示出明显的尘粒生长和结晶化证据，与棕矮星和恒星周围盘中所见的现象相似，表明这些盘具备形成岩质伴体的能力$^{[67]}$。
\subsubsection{类似行星的形成}
被抛射行星在理论上预计主要为低质量天体（$<30\,M_\oplus$，见 Ma 等人的图 1）$^{[68]}$，其平均质量取决于其宿主恒星的质量。Ma 等人$^{[68]}$ 的数值模拟表明，质量为 $1\,M_\odot$ 的恒星中，有 $17.5\%$ 会发生行星抛射事件，每颗恒星平均抛射出总质量为 $16.8\,M_\oplus$ 的物质，而单个自由漂浮行星（FFP）的典型（中位）质量约为 $0.8\,M_\oplus$。对于质量为 $0.3\,M_\odot$ 的低质量红矮星，有 $12\%$ 的恒星会发生抛射事件，每颗恒星平均抛射出总质量为 $5.1\,M_\oplus$ 的物质，而单个 FFP 的典型质量约为 $0.3\,M_\oplus$。

Hong 等人$^{[69]}$ 预测，系外卫星可以在行星—行星相互作用中被散射并成为被抛射的系外卫星。质量较高（$0.3$--$1\,M_J$）的被抛射 FFP 在理论上也是可能的，但预计十分罕见$^{[68]}$。行星的抛射可以通过行星—行星散射过程发生，或由恒星近距离掠过所触发。另一种可能性是原行星盘碎片被抛射出去，并随后演化形成行星质量天体$^{[70]}$。还有一种提出的情景是，位于倾斜的环双星轨道上的行星被抛射出去。中央双星与行星之间以及行星彼此之间的相互作用，可能导致系统中质量较低的行星被抛射$^{[71][72]}$。然而，这一机制的有效性依赖于相遇的几何构型，而该构型在观测和理论上目前均缺乏良好约束。
\subsubsection{通过年轻环星盘相互作用形成}
处于边缘引力稳定状态的年轻环星盘之间的相互作用，可能产生拉长的潮汐桥结构，这些结构可在局部发生坍缩并形成 iPMO$^{[73]}$。这些 iPMO 拥有与观测结果相符的大尺度盘结构$^{[60]}$，而被抛射行星假说难以解释这一特征。正如在梯形星团中的 iPMO 所显示的那样，它们在形成时还具有较高的多体系统比例$^{[44]}$。然而，这一机制的有效性同样依赖于相遇几何构型，而该构型在观测和理论上目前仍缺乏良好约束$^{[74]}$。
\subsubsection{其他情景}
如果恒星或棕矮星胚胎的吸积过程被中途终止，其质量可能保持在足够低的水平，从而成为行星质量天体。这种吸积中断可能发生在胚胎被抛射的情况下，或其环星盘在$O$型恒星附近经历光致蒸发时。通过“被抛射胚胎”情景形成的天体，其盘结构会较小甚至不存在，并且这类天体的双星比例会降低。也有可能，自由漂浮的行星质量天体是多种形成机制共同作用的结果$^{[70]}$。
\subsection{命运}
大多数孤立的行星质量天体将永远在星际空间中漂浮。

部分 iPMO 会与行星系统发生近距离相遇。这种罕见的相遇可能产生三种结果：iPMO 仍保持不受束缚状态；它可能与恒星形成弱引力束缚；或者它可能将一颗系外行星“踢出”系统并取而代之。数值模拟表明，绝大多数此类相遇都会导致一次俘获事件，其中 iPMO 与恒星形成弱束缚状态，其引力束缚能较低，并处在拉长的高偏心率轨道上。这类轨道并不稳定，其中约 90\% 的天体会因行星—行星相互作用而获得能量，最终再次被抛射回星际空间。只有约 1\% 的恒星会经历这种短暂的俘获过程$^{[75]}$。
\subsection{热量}
\begin{figure}[ht]
\centering
\includegraphics[width=6cm]{./figures/e7e727d8b9ebfbed.png}
\caption{一颗木星大小流浪行星的艺术构想图} \label{fig_RoPl_5}
\end{figure}
星际行星几乎不产生热量，也不会受到恒星的加热$^{[76]}$。然而，1998 年 David~J.~Stevenson 提出理论认为，一些在星际空间中漂流的行星尺度天体可能维持一层厚重的大气而不会冻结。他提出，这类大气能够依靠富含氢的大气层在高压条件下产生的远红外辐射不透明性而得以保存$^{[77]}$。

在行星系统形成过程中，若干较小的原行星天体可能会被抛射出系统$^{[78]}$。被抛射的天体将接收到更少由恒星产生的紫外辐射，而这些紫外辐射通常会剥离大气中的轻元素。即便是地球大小的天体，其引力也足以阻止大气中氢和氦的逃逸$^{[77]}$。对于地球尺度的天体而言，其核心中残余放射性同位素衰变所产生的地热能，可能使其表面温度维持在水的熔点以上$^{[77]}$，从而允许液态水海洋的存在。这类行星很可能在很长时间尺度上保持地质活动性。如果它们具有由地磁发电机产生的保护性磁层以及海底火山活动，那么热液喷口就可能为生命提供能量来源$^{[77]}$。由于其热微波辐射极其微弱，这些天体将难以被探测；不过，对于距离地球小于 1000 个天文单位的天体，其反射的太阳辐射以及远红外热辐射仍有可能被探测到$^{[79]}$。约有 $5\%$ 的地球大小被抛射行星在抛射后仍能保留其月球尺度的天然卫星。一个较大的卫星将成为显著的地质潮汐加热来源$^{[80]}$。
\subsection{列表}
下表列出了已确认或疑似的流浪行星。这些行星究竟是从围绕恒星的轨道中被抛射出去的，还是作为亚棕矮星而独立形成，目前尚不清楚。对于质量异常低的流浪行星（例如 OGLE-2012-BLG-1323 和 KMT-2019-BLG-2073），它们是否具备独立形成的能力，目前仍未知。
\subsubsection{通过直接成像发现}
这些天体是通过直接成像方法发现的。其中许多是在年轻的恒星团或恒星协会中被发现的，也已知存在少数年老的天体（例如 WISE~0855-0714）。该列表按发现年份排序。

\begin{table}[ht]
\centering
\caption{请输入表格标题}\label{tab_RoPl1}
\begin{tabular}{|c|c|c|c|c|c|c|c|}
\hline
\textbf{系外行星} & \textbf{质量（$M_J$）}& \textbf{年龄（Myr）} &\textbf{ 距离（ly）} & \textbf{光谱型} & \textbf{状态} & \textbf{恒星协会成员关系} & \textbf{发现}\\
\hline
OTS~44 & $\sim11.5$ & $0.5-3 $& 554 & M9.5 & 很可能是低质量棕矮星$^{[37]}$ & 变色龙座~I & 1998\\
\hline
S~Ori~52 & $2-8$ & $1-5 $& $1{,}150$  &  & 年龄和质量不确定；可能是前景棕矮星 & $\sigma$~猎户座星团 & 2000$^{[10]}$\\
\hline
Proplyd~061-401 & $\sim11$ & 1 & $1{,}344$ & $L4-L5$ & 候选体；本研究共发现 15 个候选体 & 猎户座星云 & 2001$^{[21]}$\\
\hline
S~Ori~70 & 3 & 3 & 1150 & T6 & 可能为混入天体？$^{[22]}$ & $\sigma$~猎户座星团 & 2002\\
\hline
Cha~110913-773444 & $5-15$ & $\sim2$ & 529 & $>$M9.5 & 已确认 & 变色龙座~I & 2004$^{[81]}$\\
\hline
SIMP~J013656.5+093347 & $11-13$ & $\sim200$ & $20-22$ & T2.5 & 候选体 & 船底座邻近运动群 & 2006$^{[82][83]}$\\
\hline
Cha~1107-7626 & $6-10$ & $1-5$ & $620$ & $L0-L1$ & 已确认 & 变色龙座~I & 2008$^{[84]}$\\
\hline
UGPS~J072227.51-054031.2 & $0.66-16.02$,$^{[85][86]}$ & 1000--5000& 13 & T9 & 质量不确定 & 无 & 2010\\
\hline
M10$-$4450 & $2-3$ & 1 & 325 & T & 候选体 & $\rho$~蛇夫座分子云 & 2010$^{[87]}$\\
\hline
WISE~1828+2650 & $3-6$ 或 $0.5-20$,$^{[88]}$ & $2-4$ 或 $0.1-10$,$^{[88]}$ & 47 & >Y2 & 候选体，可能为双星 & 无 & 2011\\
\hline
CFBDSIR~2149-0403 & $4-7$ & $110-130$ & $117-143$ & T7 & 候选体 & AB~Doradus 运动群 & 2012$^{[89]}$\\
\hline
SONYC-NGC1333-36 & $\sim6$ & 1 & 978 & L3 & 候选体；NGC~1333 中另有两个质量低于 $15\,M_J$ 的天体 & NGC~1333 & 2012$^{[90]}$\\
\hline
SSTc2d~J183037.2+011837 & $2-4$ & 3 & $848-1354$ & T? & 候选体，亦称 ID~4  & 巨蛇座核心星团$^{[91]}$（位于巨蛇座分子云） & 2012$^{[12]}$\\
\hline
PSO~J318.5-22 & $6.24-7.60$,$^{[85][86]}$ & $21-27$ & 72.32 & L7 & 已确认；亦称 2MASS~J21140802-2251358 & $\beta$~绘架座运动群 & 2013$^{[14][92]}$\\
\hline
2MASS~J2208+2921 & $11-13$  & $21-27$ & 115 & L3$\gamma$ & 候选体；需径向速度确认 & $\beta$~绘架座运动群 & 2014$^{[93]}$\\
\hline
WISE~J1741-4642 & $4-21$ & $23-130$ &  & L7pec  & 候选体  & $\beta$~绘架座或 AB~Doradus 运动群 & 2014$^{[94]}$\\
\hline
WISE~0855-0714 & $3-10$ & $>1{,}000$ & 7.1 & Y4 & 年龄不确定，但由于位于太阳邻域，推测为年老天体；$^{[95]}$ 即便在 $12$~Gyr 的高龄情形下仍为候选体（宇宙年龄为 $13.8$~Gyr）。目前已知距离最近的可能流浪行星 & 无 & 2014$^{[96]}$\\
\hline
2MASS~J12074836-3900043 & $\sim15$,$^{[97]}$ & 7--13 & 200 & L1 & 候选体；需更精确的距离测量 & TW~Hydrae 协会$^{[98]}$ & 2014$^{[99]}$\\
\hline
SIMP~J2154-1055 & 9--11 & 30--50 & 63 & L4$\beta$ & 年龄存在争议$^{[100]}$ & Argus 协会 & 2014$^{[101]}$\\
\hline
SDSS~J111010.01+011613.1 & $10.83-11.73$,$^{[85][86]}$ & $110-130$ & 63 & T5.5 & 已确认$^{[85]}$ & AB~Doradus 运动群 & 2015$^{[34]}$\\
\hline
2MASS~J11193254-1137466~AB & $4-8$ & $7-13$ & $\sim90$ & L7 & 双星候选体，其中一个分量具有候选系外卫星或可变大气$^{[55]}$ & TW~Hydrae 协会 & 2016$^{[102]}$\\
\hline
WISEA~1147 & $5-3 $& $7-13$ & $\sim100$ & L7 & 候选体 & TW~Hydrae 协会  & 2016$^{[13]}$\\
\hline
USco~J155150.2-213457 & $8-10$ & $6.907-10$ & 104 & L6 & 候选体，低表面引力 &上天蝎座协会 & 2016$^{[103]}$\\
\hline
Proplyd~133-353 & $<13$ & $0.5-1$ & $1{,}344$ & M9.5 & 具有光致蒸发盘的候选体 & 猎户座星云 & 2016$^{[60]}$\\
\hline
Cha~J11110675-7636030 & $3-6$ & $1-3$ & $520-550$  & M9-L2 & 候选体，但可能被盘包围，从而可能为亚棕矮星；该研究还提出了其他候选体 & 变色龙座~I & 2017$^{[33]}$\\
\hline
PSO~J077.1+24 & 6 & $1-2$ & 470 & L2 & 候选体；该研究还在金牛座中发表了另一候选体 & 金牛座分子云 & 2017$^{[104]}$\\
\hline
2MASS~J1115+1937 & $6^{+8}_{-4}$ & $5-45$ & 147  & L2$\gamma$ & 具有吸积盘  & 场星，可能为被抛射天体 & 2017\\
\hline
Calar~25 & $11-12$ & 120 & 435  &  & 已确认 & 昴宿星团 & 2018$^{[105]}$\\ 
\hline
2MASS~J1324+6358 & $10.7-11.8$ & $\sim150$ & $\sim33$ & T2 & 异常偏红且不太可能为双星；可靠候选体$^{[85][86]}$ & AB~Doradus 运动群 & 2007,~2018$^{[106]}$\\
\hline
WISE~J0830+2837 & $4-13$ & $>1{,}000$ & $31.3-42.7$ & $>$Y1 & 年龄不确定，但由于速度较高（高 $V_{\mathrm{tan}}$ 指示年老恒星族群）而被认为是年老天体；若年龄小于 10~Gyr 则为候选体 & 无 & 2020$^{[38]}$\\
\hline
2MASS~J0718-6415 & $3\pm1$ & $16-28$ & 30.5 & T5 & BPMG 候选成员；自转周期极短，为 $1.08$ 小时，与棕矮星 2MASS~J0348-6022 相当$^{[107][108]}$ & $\beta$~绘架座运动群 & 2021\\
\hline
DANCe~J16081299-2304316 & $3.1-6.3$ & $3-10$ & 104 & L6 & 该研究发表的至少 70 个候选体之一，光谱与 HR~8799c 相似 & 上天蝎座协会 & 2021$^{[47][50]}$\\
\hline
WISE~J2255-3118 & $2.15-2.59$ & 24 & $\sim45$ & T8 & 非常偏红，候选体$^{[85][86]}$，已确认？$^{[109]}$ & $\beta$~绘架座运动群 & 2011,~2021$^{[49]}$\\
\hline
WISE~J024124.73-365328.0 & $4.64-5.30$ & 45 & $\sim61$ & T7 & 候选体$^{[85][86]}$ & Argus 协会 & 2012,~2021$^{[49]}$\\
\hline
2MASS~J0013-1143 & $7.29-8.25$ & 45 & $\sim82$ & T4 & 双星候选体或复合大气；候选体$^{[85][86]}$ & Argus 协会 & 2017,~2021$^{[49]}$\\
\hline
SDSS~J020742.48+000056.2 & $7.11-8.61$ & 45 & $\sim112$ & T4.5 & 候选体$^{[85][86]}$ & Argus 协会 & 2002,~2021$^{[49]}$\\
\hline
2MASSI~J0453264-175154 & $12.68-12.98$ & 24 & $\sim99$ & L2.5$\beta$ & 低表面引力，候选体$^{[85][86]}$ & $\beta$~绘架座运动群 & 2003,~2023$^{[85][86]}$\\
\hline
CWISE~J0506+0738 & $7\pm2$ & 22 & 104 & $L8\gamma-T0\gamma$ & BPMG 候选成员；近红外颜色极端偏红$^{[110]}$ & $\beta$~绘架座运动群 & 2023\\
\hline
\end{tabular}
\end{table}
\subsubsection{通过引力微透镜发现}
这些天体是通过引力微透镜方法发现的。通过引力微透镜发现的流浪行星只能依赖透镜事件本身进行研究。其中一些天体也可能是在一颗未被观测到的恒星周围处于宽轨道上的系外行星$^{[111]}$。
\begin{table}[ht]
\centering
\caption{请输入表格标题}\label{tab_RoPl2}
\begin{tabular}{|c|c|c|c|c|c|}
\hline
\textbf{系外行星} & \textbf{质量（$M_J$}） & \textbf{质量（$M_\oplus$）}  & \textbf{距离（ly}） & \textbf{状态}  & \textbf{发现} \\
\hline
KMT-2023-BLG-2669 & $0.025-0.25$ & $8-80$ & & 候选体；需确定距离 & 2024$^{[112]}$\\
\hline
OGLE-2012-BLG-1323 & $0.0072-0.072$ & $2.3-23$ & &  候选体；需确定距离 & 2017$^{[113][114][115][116]}$\\
\hline
OGLE-2017-BLG-0560 & $1.9-20$ & $604-3{,}256$ &  & 候选体；需确定距离 & 2017$^{[114][115][116]}$\\
\hline
MOA-2015-BLG-337L & 9.85 & $3{,}130$ & $23{,}156$ & 可能是一个棕矮星双星系统 & 2018$^{[117][118]}$\\
\hline
KMT-2019-BLG-2073 & 0.19 & 59 &  & 候选体；需确定距离 & 2020$^{[119]}$\\
\hline
OGLE-2016-BLG-1928 & $0.001-0.006$ & $0.3-2$ & $30{,}000-180{,}000$  & 候选体 & 2020$^{[111]}$\\
\hline
OGLE-2019-BLG-0551 & $0.0242-0.3$ & $7.69-95$ &  &  刻画程度较低$^{[120]}$ & 2020$^{[120]}$\\
\hline
VVV-2012-BLG-0472L & 10.5 & $3{,}337$ & $3{,}200$ &  & 2022$^{[121]}$\\
\hline
MOA-9y-770L & 0.07 & $22.3^{+42.2}_{-17.4}$ & $22{,}700$ & & 2023$^{[122]}$\\
\hline
MOA-9y-5919L & 0.0012 或 0.0024 & $0.37^{+1.11}_{-0.27}$ 或 $0.75^{+1.23}_{-0.46}$ & $14{,}700$ 或 $19{,}300$ &  & 2023$^{[122]}$\\
\hline
OGLE-2017-BLG-1170L & $3.06^{+1.34}_{-1.16}$ &  & $24{,}700$ & 候选体 & 2019$^{[123]}$\\
\hline
OGLE-2017-BLG-1170L & $1.85^{+0.79}_{-0.70}$ & &  $24{,}700$ & 候选体  & 2019$^{[123]}$\\
\hline
\end{tabular}
\end{table}
\subsubsection{通过凌星法发现}
\begin{table}[ht]
\centering
\caption{请输入表格标题}\label{tab_RoPl3}
\begin{tabular}{|c|c|c|c|c|c|c|c|}
\hline
\textbf{系外行星} & \textbf{ 质量（$M_J$）} & \textbf{年龄（Myr）} &\textbf{ 距离（ly）} & \textbf{光谱型} & \textbf{状态} & \textbf{ 恒星协会成员关系} &\textbf{发现} \\
\hline
J1407b  & $<6$ &   & $<451$ &  & 候选的 ALMA 探测体；尽管该天体的亮度和位置与 2007 年遮掩半人马座 V1400 恒星的天体一致，但仍需要 ALMA 的后续观测以确认其是否在运动，更不用说是否朝着正确的方向运动$^{[124]}$  & 无  & 2012,~2020$^{[124]}$\\
\hline
\end{tabular}
\end{table}
\subsection{另见}
\begin{itemize}
\item Capotauro —— 位于牧夫座中的未确认天文天体  
\item 星系际恒星（Intergalactic star）—— 不受任何星系引力束缚的恒星  
\item 星际天体（Interstellar object）—— 位于星际空间、且不受任何恒星引力束缚的天文天体  
\item ʻOumuamua —— 于 2017 年穿越太阳系的星际天体  
\item 流浪黑洞（Rogue black hole）—— 引力上不受束缚的黑洞  
\item 流浪的河外行星（Rogue extragalactic planets）—— 位于银河系之外的流浪行星  
\item 潮汐剥离的系外卫星（Tidally detached exomoon）—— 原本为卫星、后成为流浪行星的天体
\end{itemize}  
\subsubsection{虚构作品中}
\begin{itemize}
\item 《一桶空气》（A~Pail~of~Air，1951）—— Fritz~Leiber 创作的科幻短篇小说，其中地球被黑洞拉出太阳系。尽管文中明确指出地球仍绕黑洞运行，但其整体效果等同于地球作为流浪行星被抛射出太阳系$^{[125][126]}$  
\item 《第九行星的秘密》（The~Secret~of~the~Ninth~Planet，1959）—— Donald~A.~Wollheim 的小说，其中冥王星被揭示为一颗被俘获的流浪行星，其上仍存有残存文明$^{[127]}$  
\item 《太空：1999》（Space:~1999，1975--77）—— 英国科幻电视剧，剧情设定中月球因一次热核爆炸而被抛射出太阳系$^{[125]}$  
\item 《流浪行星》（Rogue~Planet，2002）—— 《星际迷航：进取号》中的一集  
\item 《密特罗德 Prime 2：回音》（Metroid~Prime~2:~Echoes，2004）—— 设定于一颗流浪行星上的电子游戏，其生态系统由被当地居民称为“以太之光”（Light~of~Aether）的自然能量维系  
\item 《雷米娜》（Remina，2004--2005）—— 伊藤润二创作的恐怖漫画，描绘了一颗具有意识、能够吞噬行星和恒星的流浪行星  
\item 《忧郁症》（Melancholia，2011）—— Lars~von~Trier 执导的科幻电影，讲述一颗流浪行星与地球即将发生碰撞的故事  
\item 《黑暗伊甸园》（Dark~Eden，2012）—— Chris~Beckett 创作的社会科幻小说  
\item 《流浪地球》（The~Wandering~Earth，2019）—— 郭帆执导的科幻电影，讲述人类将地球从太阳系人工迁移至半人马座 $\alpha$ 星系的故事$^{[125]}$  
\item 《双子座家庭娱乐》（Gemini~Home~Entertainment，2019--）—— Remy~Abode 创作的恐怖选集网络剧，其主要反派是一颗名为“The~Iris”的具备意识的流浪行星，策划了对太阳系，尤其是地球与海王星的入侵$^{[128]}$  
\item 《卡罗尔与世界的终结》（Carol~\&~the~End~of~the~World，2023）—— Dan~Guterman 创作的成人向动画喜剧迷你剧
\end{itemize}
\subsection{参考文献}
\begin{enumerate}
\item Shostak, Seth（2005 年 2 月 24 日）。“孤儿行星：艰难的人生”。Space.com。检索于 2020 年 11 月 13 日。  
\item Lloyd, Robin（2001 年 4 月 18 日）。“自由漂浮行星——英国团队重新提出存疑主张”。Space.com。原文存档于 2008 年 10 月 13 日。  
\item “孤儿‘行星’的发现受到新模型挑战”。NASA 天体生物学项目。2001 年 4 月 18 日。原文存档于 2009 年 3 月 22 日。  
\item Kirkpatrick, J. Davy; Gelino, Christopher R.; Faherty, Jacqueline K.; Meisner, Aaron M.; Caselden, Dan; Schneider, Adam C.; Marocco, Federico; Cayago, Alfred J.; Smart, R. L.; Eisenhardt, Peter R.; Kuchner, Marc J.; Wright, Edward L.; Cushing, Michael C.; Allers, Katelyn N.; Bardalez Gagliuffi, Daniella C.（2021 年 3 月 1 日）。“基于 20 pc 全天区 525 颗 L、T 与 Y 矮星普查的场亚恒星质量函数”。《天体物理学杂志增刊系列》，253（1）：7。arXiv:2011.11616。Bibcode:2021ApJS..253....7K。doi:10.3847/1538-4365/abd107。ISSN 0067-0049。  
\item Neil deGrasse Tyson 在《宇宙：时空奥德赛》中提出的观点，引自《国家地理》。  
\item “研究团队发现，该任务将提供比当前估计至少精确 10 倍的流浪行星数量统计，目前银河系中的估计范围从数百亿到数万亿不等。”SciTechDaily。  
\item “我们应该多担心流浪行星撞击地球？”。2024 年 4 月 6 日。  
\item 系外行星工作组——《“行星”定义立场声明》（国际天文学联合会，IAU）。原文于 2006 年 9 月 16 日存档于 Wayback Machine。  
\item “流浪行星的发现促使天文学家重新思考理论”。CNN。原文存档于 2011 年 6 月 28 日。  
\item Zapatero Osorio, M. R.（2000 年 10 月 6 日）。“在 $\sigma$~猎户座星团中发现年轻的孤立行星质量天体”。《科学》，290（5489）：103--107。Bibcode:2000Sci...290..103Z。doi:10.1126/science.290.5489.103。PMID 11021788。  
\item Lucas, P. W.; Roche, P. F.（2000 年 6 月 1 日）。“猎户座中一群极年轻的棕矮星与自由漂浮行星”。《皇家天文学会月刊》，314（4）：858--864。arXiv:astro-ph/0003061。Bibcode:2000MNRAS.314..858L。doi:10.1046/j.1365-8711.2000.03515.x。ISSN 0035-8711。S2CID 119002349。  
\item Spezzi, L.; Alves de Oliveira, C.; Moraux, E.; Bouvier, J.; Winston, E.; Hudelot, P.; Bouy, H.; Cuillandre, J.-C.（2012 年 9 月 1 日）。“在巨蛇座核心区寻找行星质量 T 矮星”。《天文学与天体物理学》，545：A105。arXiv:1208.0702。Bibcode:2012A\&A...545A.105S。doi:10.1051/0004-6361/201219559。ISSN 0004-6361。S2CID 119232214。  
\item Schneider, Adam C.（2016 年 4 月 21 日）。“WISEA~J114724.10$-$204021.3：TW~Hya 协会中的一颗自由漂浮行星质量成员”。《天体物理学杂志快报》，822（1）：L1。arXiv:1603.07985。Bibcode:2016ApJ...822L...1S。doi:10.3847/2041-8205/822/1/L1。S2CID 30068452。  
\item Liu, Michael C.（2013 年 11 月 10 日）。“极端偏红的年轻 L 矮星 PSO~J318.5338$-$22.8603：直接成像年轻气态巨行星的自由漂浮行星质量类比体”。《天体物理学杂志快报》，777（1）：L20。arXiv:1310.0457。Bibcode:2013ApJ...777L..20L。doi:10.1088/2041-8205/777/2/L20。S2CID 54007072。  
\item Bennett, D. P.; Batista, V.; 等（2013 年 12 月 13 日）。“一颗绕气态巨行星运行的亚地球质量卫星，或银河系核球中的高速行星系统”。《天体物理学杂志》，785（2）：155。arXiv:1312.3951。Bibcode:2014ApJ...785..155B。doi:10.1088/0004-637X/785/2/155。S2CID 118327512。  
\item Mróz, Przemek; 等（2020 年）。“在持续时间最短的引力微透镜事件中探测到一颗地球质量的流浪行星候选体”。《天体物理学杂志快报》，903（1）：L11。arXiv:2009.12377。Bibcode:2020ApJ...903L..11M。doi:10.3847/2041-8213/abbfad。  
\item “ESO 望远镜帮助揭示迄今最大规模的流浪行星群”。欧洲南方天文台。2021 年 12 月 22 日。  
\item “数十亿颗无恒星行星潜伏于黑暗云团的摇篮中”。日本国家天文台（NAOJ）。2021 年 12 月 23 日。  
\item Shen, Zili（2021 年 12 月 30 日）。“游荡的行星”。Astrobites。  
\item “银河系中发现最大规模的自由漂浮行星集合——KPNO”。kpno.noirlab.edu。  
\item Lucas, P. W.; Roche, P. F.; Allard, France; Hauschildt, P. H.（2001 年 9 月 1 日）。“猎户座中亚恒星天体的红外光谱研究”。《皇家天文学会月刊》，326（2）：695--721。arXiv:astro-ph/0105154。Bibcode:2001MNRAS.326..695L。doi:10.1046/j.1365-8711.2001.04666.x。ISSN 0035-8711。S2CID 280663。  
\item Caballero, José A.（2018 年 9 月 1 日）。“低于氘燃烧质量极限的亚恒星天体综述：行星、棕矮星，还是其他？”《地球科学》，8（10）：362。arXiv:1808.07798。Bibcode:2018Geosc...8..362C。doi:10.3390/geosciences8100362。  
\item Oasa, Yumiko; Tamura, Motohide; Sugitani, Koji（1999 年 11 月 1 日）。“变色龙座 I 暗云核心的深度近红外巡天”。《天体物理学杂志》，526（1）：336--343。Bibcode:1999ApJ...526..336O。doi:10.1086/307964。ISSN 0004-637X。S2CID 120597899。  
\item Luhman, K. L.; Peterson, Dawn E.; Megeath, S. T.（2004 年 12 月 1 日）。“对变色龙座中已知质量最低棕矮星的光谱确认”。《天体物理学杂志》，617（1）：565--568。arXiv:astro-ph/0411445。Bibcode:2004ApJ...617..565L。doi:10.1086/425228。ISSN 0004-637X。S2CID 18157277。  
\item Cartwright, Jon（2011 年 5 月 18 日）。“无家可归的行星在银河系中或许很常见”。Science~Now。  
\item “没有恒星的行星：发现了一类新的行星”。Physorg.com，2011 年 5 月 18 日。  
\item Sumi, T.; 等（2011 年）。“通过引力微透镜探测到不受束缚或远轨道的行星质量族群”。《自然》，473（7347）：349--352。arXiv:1105.3544。Bibcode:2011Natur.473..349S。doi:10.1038/nature10092。PMID 21593867。S2CID 4422627。  
\item “研究人员称银河系可能充满‘游牧行星’”。斯坦福大学。2012 年 2 月 23 日。  
\item Mróz, P.; 等（2017 年）。“不存在大量不受束缚或宽轨道的木星质量行星族群”。《自然》，548（7666）：183--186。arXiv:1707.07634。Bibcode:2017Natur.548..183M。doi:10.1038/nature23276。PMID 28738410。S2CID 4459776。  
\item Gough, Evan（2020 年 10 月 1 日）。“在银河系中发现一颗不依附于任何恒星、自由漂浮的地球质量流浪行星”。Universe~Today。
\item Redd, Nola Taylor（2020 年 10 月 19 日）。“在银河系中发现漂流的流浪岩质行星——这一微小世界及类似天体或可帮助天文学家探究行星形成之谜”。《科学美国人》。检索于 2020 年 10 月 19 日。  
\item Saumon, D.; Marley, Mark S.（2008 年 12 月 1 日）。“L 与 T 矮星在颜色—星等图中的演化”。《天体物理学杂志》，689（2）：1327--1344。arXiv:0808.2611。Bibcode:2008ApJ...689.1327S。doi:10.1086/592734。ISSN 0004-637X。S2CID 15981010。  
\item Esplin, T. L.; Luhman, K. L.; Faherty, J. K.; Mamajek, E. E.; Bochanski, J. J.（2017 年 8 月 1 日）。“变色龙座 I 恒星形成区中行星质量棕矮星的巡天”。《天文学杂志》，154（2）：46。arXiv:1706.00058。Bibcode:2017AJ....154...46E。doi:10.3847/1538-3881/aa74e2。ISSN 0004-6256。  
\item Gagné, Jonathan（2015 年 7 月 20 日）。“SDSS~J111010.01+011613.1：AB~Doradus 运动群中的一颗新的行星质量 T 矮星成员”。《天体物理学杂志快报》，808（1）：L20。arXiv:1506.04195。Bibcode:2015ApJ...808L..20G。doi:10.1088/2041-8205/808/1/L20。S2CID 118834638。  
\item Leggett, S. K.; Tremblin, P.; Esplin, T. L.; Luhman, K. L.; Morley, Caroline V.（2017 年 6 月 1 日）。“Y 型棕矮星：基于新的天体测量、统一化测光与近红外光谱的质量与年龄估计”。《天体物理学杂志》，842（2）：118。arXiv:1704.03573。Bibcode:2017ApJ...842..118L。doi:10.3847/1538-4357/aa6fb5。ISSN 0004-637X。  
\item Luhman, Kevin L.; Esplin, Taran L.（2016 年 9 月）。“已知最冷棕矮星的光谱能量分布”。《天文学杂志》，152（2）：78。arXiv:1605.06655。Bibcode:2016AJ....152...78L。doi:10.3847/0004-6256/152/3/78。S2CID 118577918。  
\item Luhman, Kevin L.（2005 年 2 月 10 日）。“斯皮策对已知质量最低、具有环星盘棕矮星的识别”。《天体物理学杂志快报》，620（1）：L51--L54。arXiv:astro-ph/0502100。Bibcode:2005ApJ...620L..51L。doi:10.1086/428613。S2CID 15340083。  
\item Bardalez Gagliuffi, Daniella C.; Faherty, Jacqueline K.; Schneider, Adam C.; Meisner, Aaron; Caselden, Dan; Colin, Guillaume; Goodman, Sam; Kirkpatrick, J. Davy; Kuchner, Marc; Gagné, Jonathan; Logsdon, Sarah E.（2020 年 6 月 1 日）。“WISEA~J083011.95+283716.0：缺失环节的行星质量天体”。《天体物理学杂志》，895（2）：145。arXiv:2004.12829。Bibcode:2020ApJ...895..145B。doi:10.3847/1538-4357/ab8d25。S2CID 216553879。  
\item Kirkpatrick, J. Davy; Marocco, Federico; Caselden, Dan; Meisner, Aaron M.; Faherty, Jacqueline K.; Schneider, Adam C.; Kuchner, Marc J.; Casewell, S. L.; Gelino, Christopher R.; Cushing, Michael C.; Eisenhardt, Peter R.; Wright, Edward L.; Schurr, Steven D.（2021 年 7 月 1 日）。“谜一样的棕矮星 WISEA~J153429.75$-$104303.3（又名‘The Accident’）”。《天体物理学杂志》，915（1）：L6。arXiv:2106.13408。Bibcode:2021ApJ...915L...6K。doi:10.3847/2041-8213/ac0437。ISSN 0004-637X。  
\item Joergens, V.; Bonnefoy, M.; Liu, Y.; Bayo, A.; Wolf, S.; Chauvin, G.; Rojo, P.（2013 年）。“OTS~44：位于行星边界处的盘与吸积”。《天文学与天体物理学》，558（7）：L7。arXiv:1310.1936。Bibcode:2013A\&A...558L...7J。doi:10.1051/0004-6361/201322432。S2CID 118456052。  
\item Dupuy, Trent J.; Liu, Michael C.（2012 年 8 月 1 日）。“夏威夷红外视差计划 I：超冷双星与 L/T 过渡”。《天体物理学杂志增刊系列》，201（2）：19。arXiv:1201.2465。Bibcode:2012ApJS..201...19D。doi:10.1088/0067-0049/201/2/19。ISSN 0067-0049。  
\item Zhang, Zhoujian; Liu, Michael C.; Best, William M. J.; Dupuy, Trent J.; Siverd, Robert J.（2021 年 4 月 1 日）。“夏威夷红外视差计划 V：附近年轻运动群中的新 T 矮星成员与候选成员”。《天体物理学杂志》，911（1）：7。arXiv:2102.05045。Bibcode:2021ApJ...911....7Z。doi:10.3847/1538-4357/abe3fa。ISSN 0004-637X。  
\item Langeveld, Adam B.; Scholz, Aleks; Mužić, Koraljka; Jayawardhana, Ray; Capela, Daniel; Albert, Loïc; Doyon, René; Flagg, Laura; de Furio, Matthew; Johnstone, Doug; Lafrèniere, David; Meyer, Michael（2024 年 10 月 1 日）。“JWST/NIRISS 面向年轻棕矮星与自由漂浮行星的深度光谱巡天”。《天文学杂志》，168（4）：179。arXiv:2408.12639。Bibcode:2024AJ....168..179L。doi:10.3847/1538-3881/ad6f0c。ISSN 0004-6256。  
\item Pearson, Samuel G.; McCaughrean, Mark J.（2023 年 10 月 2 日）。“梯形星团中的木星质量双星天体”。第 24 页。arXiv:2310.01231 [astro-ph.EP]。  
\item Portegies Zwart, Simon; Hochart, Erwan（2024 年 7 月 2 日）。“年轻恒星团中宽分离木星质量双星天体的起源与演化”。SciPost，3（1）：19。arXiv:2312.04645。Bibcode:2024ScPA....3....1P。doi:10.21468/SciPostAstro.3.1.001。  
\item Luhman, K. L.（2024 年 10 月 14 日）。“基于 JWST/NIRCam 的猎户座星云星团亚恒星成员候选体”。《天文学杂志》，168（6）：230。arXiv:2410.10406。Bibcode:2024AJ....168..230L。doi:10.3847/1538-3881/ad812a。  
\item Miret-Roig, Núria; Bouy, Hervé; Raymond, Sean N.; Tamura, Motohide; Bertin, Emmanuel; Barrado, David; Olivares, Javier; Galli, Phillip A. B.; Cuillandre, Jean-Charles; Sarro, Luis Manuel; Berihuete, Angel（2021 年 12 月 22 日）。“上天蝎座年轻恒星协会中丰富的自由漂浮行星族群”。Nature Astronomy，6：89--97。arXiv:2112.11999。Bibcode:2022NatAs...6...89M。doi:10.1038/s41550-021-01513-x。ISSN 2397-3366。S2CID 245385321。另见 Nature SharedIt 文章链接；ESO 文章链接。  
\item Béjar, V. J. S.; Martín, Eduardo L.（2018 年 1 月 1 日）。《年轻恒星团中的棕矮星与自由漂浮行星》。Bibcode:2018haex.bookE..92B。  
\item Zhang, Zhoujian; Liu, Michael C.; Best, William M. J.; Dupuy, Trent J.; Siverd, Robert J.（2021 年 4 月 1 日）。“夏威夷红外视差计划 V：附近年轻运动群中的新 T 矮星成员与候选成员”。《天体物理学杂志》，911（1）：7。arXiv:2102.05045。Bibcode:2021ApJ...911....7Z。doi:10.3847/1538-4357/abe3fa。ISSN 0004-637X。  
\item Bouy, H.; Tamura, M.; Barrado, D.; Motohara, K.; Castro Rodríguez, N.; Miret-Roig, N.; Konishi, M.; Koyama, S.; Takahashi, H.; Huélamo, N.; Bertin, E.; Olivares, J.; Sarro, L. M.; Berihuete, A.; Cuillandre, J.-C.（2022 年 8 月 1 日）。“上天蝎座与蛇夫座自由漂浮行星候选体的红外光谱研究”。《天文学与天体物理学》，664：A111。arXiv:2206.00916。Bibcode:2022A\&A...664A.111B。doi:10.1051/0004-6361/202243850。ISSN 0004-6361。S2CID 249282287。  
\item Raymond, Sean; Bouy, Núria Miret-Roig \& Hervé（2021 年 12 月 22 日）。“我们发现了一组怪物般巨大的气态巨行星流浪者画廊”。Nautilus。检索于 2021 年 12 月 23 日。  
\item Pandey, Shreesham; Pandey, Diva（2025 年）。“流浪行星塑造系外行星总体特征：动力学、生化与动物学影响”。2025 年 Sagan 夏季研讨会，NExScI（NASA 系外行星科学研究所），加州理工学院，美国加利福尼亚州帕萨迪纳。doi:10.13140/RG.2.2.23859.28964。  
\item Boss, Alan P.（2001 年 4 月 1 日）。“通过原恒星坍缩与碎裂形成行星质量天体”。《天体物理学杂志》，551（2）：L167--L170。Bibcode:2001ApJ...551L.167B。doi:10.1086/320033。ISSN 0004-637X。S2CID 121261733。  
\item Gahm, G. F.; Grenman, T.; Fredriksson, S.; Kristen, H.（2007 年 4 月 1 日）。“作为棕矮星与自由漂浮行星质量天体种子的‘小球状体’”。《天文学杂志》，133（4）：1795--1809。Bibcode:2007AJ....133.1795G。doi:10.1086/512036。ISSN 0004-6256。S2CID 120588285。  
\item Limbach, Mary Anne; Vos, Johanna M.; Winn, Joshua N.; Heller, René; Mason, Jeffrey C.; Schneider, Adam C.; Dai, Fei（2021 年 9 月 1 日）。“关于凌星孤立行星质量天体之系外卫星的探测”。《天体物理学杂志》，918（2）：L25。arXiv:2108.08323。Bibcode:2021ApJ...918L..25L。doi:10.3847/2041-8213/ac1e2d。ISSN 0004-637X。  
\item Luhman, K. L.; Adame, Lucía; D'Alessio, Paola; Calvet, Nuria; Hartmann, Lee; Megeath, S. T.; Fazio, G. G.（2005 年 12 月 1 日）。“发现一颗具有环星盘的行星质量棕矮星”。《天体物理学杂志》，635（1）：L93--L96。arXiv:astro-ph/0511807。Bibcode:2005ApJ...635L..93L。doi:10.1086/498868。ISSN 0004-637X。  
\item Jayawardhana, Ray; Ivanov, Valentin D.（2006 年 8 月 1 日）。“对具有盘的年轻行星质量候选体的光谱观测”。《天体物理学杂志》，647（2）：L167--L170。arXiv:astro-ph/0607152。Bibcode:2006ApJ...647L.167J。doi:10.1086/507522。ISSN 0004-637X。  
\item Rilinger, Anneliese M.; Espaillat, Catherine C.（2021 年 11 月 1 日）。“棕矮星周围原行星盘的盘质量与尘埃演化”。《天体物理学杂志》，921（2）：182。arXiv:2106.05247。Bibcode:2021ApJ...921..182R。doi:10.3847/1538-4357/ac09e5。ISSN 0004-637X。  
\item Zapatero Osorio, M. R.; Caballero, J. A.; Béjar, V. J. S.; Rebolo, R.; Barrado Y Navascués, D.; Bihain, G.; Eislöffel, J.; Martín, E. L.; Bailer-Jones, C. A. L.; Mundt, R.; Forveille, T.; Bouy, H.（2007 年 9 月 1 日）。“$\sigma$~猎户座中行星质量天体的盘”。《天文学与天体物理学》，472（1）：L9--L12。Bibcode:2007A\&A...472L...9Z。doi:10.1051/0004-6361:20078116。ISSN 0004-6361。  
\item Fang, Min; Kim, Jinyoung Serena; Pascucci, Ilaria; Apai, Dániel; Manara, Carlo Felice（2016 年 12 月 1 日）。“猎户座中具有光致蒸发盘的行星质量天体候选体”。《天体物理学杂志快报》，833（2）：L16。arXiv:1611.09761。Bibcode:2016ApJ...833L..16F。doi:10.3847/2041-8213/833/2/L16。ISSN 0004-637X。  
\item Best, William M. J.; Liu, Michael C.; Magnier, Eugene A.; Bowler, Brendan P.; Aller, Kimberly M.; Zhang, Zhoujian; Kotson, Michael C.; Burgett, W. S.; Chambers, K. C.; Draper, P. W.; Flewelling, H.; Hodapp, K. W.; Kaiser, N.; Metcalfe, N.; Wainscoat, R. J.（2017 年 3 月 1 日）。“利用 Pan-STARRS1 与 WISE 搜寻 L/T 过渡矮星 III：金牛座与天蝎—半人马座中的年轻 L 矮星发现与自行星表”。《天体物理学杂志》，837（1）：95。arXiv:1702.00789。Bibcode:2017ApJ...837...95B。doi:10.3847/1538-4357/aa5df0。hdl:1721.1/109753。ISSN 0004-637X。  
\item Scholz, Aleks; Muzic, Koraljka; Jayawardhana, Ray; Almendros-Abad, Victor; Wilson, Isaac（2023 年 5 月 1 日）。“年轻行星质量天体周围的盘：NGC~1333 的超深斯皮策成像”。《天文学杂志》，165（5）：196。arXiv:2303.12451。Bibcode:2023AJ....165..196S。doi:10.3847/1538-3881/acc65d。hdl:10023/27429。ISSN 0004-6256。  
\item Alves de Oliveira, C.; Moraux, E.; Bouvier, J.; Duchêne, G.; Bouy, H.; Maschberger, T.; Hudelot, P.（2013 年 1 月 1 日）。“IC~348 中棕矮星候选体的光谱观测及其亚恒星初始质量函数在行星质量范围内的确定”。《天文学与天体物理学》，549：A123。arXiv:1211.4029。Bibcode:2013A\&A...549A.123A。doi:10.1051/0004-6361/201220229。ISSN 0004-6361。  
\item Boucher, Anne; Lafrenière, David; Gagné, Jonathan; Malo, Lison; Faherty, Jacqueline K.; Doyon, René; Chen, Christine H.（2016 年 11 月 1 日）。“BANYAN VIII：具有候选环星盘的新低质量恒星与棕矮星”。《天体物理学杂志》，832（1）：50。arXiv:1608.08259。Bibcode:2016ApJ...832...50B。doi:10.3847/0004-637X/832/1/50。ISSN 0004-637X。  
\item Theissen, Christopher A.; Burgasser, Adam J.; Bardalez Gagliuffi, Daniella C.; Hardegree-Ullman, Kevin K.; Gagné, Jonathan; Schmidt, Sarah J.; West, Andrew A.（2018 年 1 月 1 日）。“2MASS~J11151597+1937266：一颗年轻、多尘、孤立的行星质量天体，可能具有一颗宽分离的恒星伴星”。《天体物理学杂志》，853（1）：75。arXiv:1712.03964。Bibcode:2018ApJ...853...75T。doi:10.3847/1538-4357/aaa0cf。ISSN 0004-637X。
\item Flagg, Laura; Scholz, Aleks; Almendros-Abad, V.; Jayawardhana, Ray; Damian, Belinda; Muzic, Koraljka; Natta, Antonella; Pinilla, Paola; Testi, Leonardo（2025 年）。“在一颗正在主动吸积的行星质量天体周围的盘中探测到烃类物质”。《天体物理学杂志》，986（2）：200。arXiv:2505.13714。Bibcode:2025ApJ...986..200F。doi:10.3847/1538-4357/add71d。  
\item Damian, Belinda; Scholz, Aleks; Jayawardhana, Ray; Almendros-Abad, V.; Flagg, Laura; Mužić, Koraljka; Natta, Antonella; Pinilla, Paola; Testi, Leonardo（2025 年）。“利用 JWST 对自由漂浮行星质量天体及其盘的光谱观测”。《天文学杂志》，170（2）：127。arXiv:2507.05155。Bibcode:2025AJ....170..127D。doi:10.3847/1538-3881/adea50。  
\item Ma, Sizheng; Mao, Shude; Ida, Shigeru; Zhu, Wei; Lin, Douglas N. C.（2016 年 9 月 1 日）。“基于核吸积理论的自由漂浮行星：引力微透镜预测”。《皇家天文学会月刊》，461（1）：L107--L111。arXiv:1605.08556。Bibcode:2016MNRAS.461L.107M。doi:10.1093/mnrasl/slw110。ISSN 0035-8711。  
\item Hong, Yu-Cian; Raymond, Sean N.; Nicholson, Philip D.; Lunine, Jonathan I.（2018 年 1 月 1 日）。“无辜的旁观者：行星—行星散射过程中系外卫星的轨道动力学”。《天体物理学杂志》，852（2）：85。arXiv:1712.06500。Bibcode:2018ApJ...852...85H。doi:10.3847/1538-4357/aaa0db。ISSN 0004-637X。  
\item Miret-Roig, Núria（2023 年 3 月 1 日）。“自由漂浮行星的起源”。《天体物理与空间科学》，368（3）：17。arXiv:2303.05522。Bibcode:2023Ap\&SS.368...17M。doi:10.1007/s10509-023-04175-5。ISSN 0004-640X。  
\item Chen, Cheng; Martin, Rebecca G.; Lubow, Stephen H.; Nixon, C. J.（2024 年 1 月 1 日）。“倾斜的环双星行星系统作为高效自由漂浮行星前身”。《天体物理学杂志》，961（1）：L5。arXiv:2310.15603。Bibcode:2024ApJ...961L...5C。doi:10.3847/2041-8213/ad17c5。ISSN 0004-637X。  
\item Portegies Zwart, Simon; Hochart, Erwan（2024 年 7 月 2 日）。“年轻恒星团中宽分离木星质量双星天体的起源与演化”。SciPost Astronomy，3（1）：001。arXiv:2312.04645。Bibcode:2024ScPA....3....1P。doi:10.21468/SciPostAstro.3.1.001。  
\item “PMO”。www.news.uzh.ch。2025 年 2 月 27 日。  
\item Fu, Zhihao; Deng, Hongping; Lin, Douglas N. C.; Mayer, Lucio（2025 年 2 月 28 日）。“通过环星盘相互作用形成自由漂浮行星质量天体”。《科学进展》，11（9）：eadu6058。arXiv:2410.21180。Bibcode:2025SciA...11.6058F。doi:10.1126/sciadv.adu6058。PMC 11864182。PMID 40009673。  
\item Goulinski, Nadav; Ribak, Erez N.（2018 年 1 月 1 日）。“行星系统对自由漂浮行星的俘获”。《皇家天文学会月刊》，473（2）：1589--1595。arXiv:1705.10332。Bibcode:2018MNRAS.473.1589G。doi:10.1093/mnras/stx2506。ISSN 0035-8711。  
\item Raymond, Sean（2005 年 4 月 9 日）。“黑暗中的生命”。Aeon。检索于 2016 年 4 月 9 日。  
\item Stevenson, David J.; Stevens, C. F.（1999 年）。“星际空间中可维持生命的行星？”。《自然》，400（6739）：32。Bibcode:1999Natur.400...32S。doi:10.1038/21811。PMID 10403246。S2CID 4307897。  
\item Lissauer, J. J.（1987 年）。“行星吸积的时间尺度与原行星盘结构”。《伊卡洛斯》，69（2）：249--265。Bibcode:1987Icar...69..249L。doi:10.1016/0019-1035(87)90104-7。hdl:2060/19870013947。  
\item Abbot, Dorian S.; Switzer, Eric R.（2011 年 6 月 2 日）。“Steppenwolf：关于星际空间中宜居行星的一个提案”。《天体物理学杂志》，735（2）：L27。arXiv:1102.1108。Bibcode:2011ApJ...735L..27A。doi:10.1088/2041-8205/735/2/L27。S2CID 73631942。  
\item Debes, John H.; Steinn Sigurðsson（2007 年 10 月 20 日）。“具有卫星的被抛射类地行星的存活率”。《天体物理学杂志快报》，668（2）：L167--L170。arXiv:0709.0945。Bibcode:2007ApJ...668L.167D。doi:10.1086/523103。S2CID 15782213。  
\item Luhman, Kevin L.（2005 年 12 月 10 日）。“发现一颗具有环星盘的行星质量棕矮星”。《天体物理学杂志快报》，635（1）：L93--L96。arXiv:astro-ph/0511807。Bibcode:2005ApJ...635L..93L。doi:10.1086/498868。S2CID 11685964。  
\item Artigau, Étienne; Doyon, René; Lafrenière, David; Nadeau, Daniel; Robert, Jasmin; Albert, Loïc（无日期）。“发现北半球最亮的 T 矮星”。《天体物理学杂志快报》，651（1）：L57。arXiv:astro-ph/0609419。Bibcode:2006ApJ...651L..57A。doi:10.1086/509146。ISSN 1538-4357。S2CID 118943169。  
\item Gagné, Jonathan; Faherty, Jacqueline K.; Burgasser, Adam J.; Artigau, Étienne; Bouchard, Sandie; Albert, Loïc; Lafrenière, David; Doyon, René; Bardalez-Gagliuffi, Daniella C.（2017 年 5 月 15 日）。“SIMP~J013656.5+093347 很可能是船底座邻近运动群中的行星质量天体”。《天体物理学杂志》，841（1）：L1。arXiv:1705.01625。Bibcode:2017ApJ...841L...1G。doi:10.3847/2041-8213/aa70e2。ISSN 2041-8213。S2CID 119024210。  
\item Luhman, K. L.; Allen, L. E.; Allen, P. R.; Gutermuth, R. A.; Hartmann, L.; Mamajek, E. E.; Megeath, S. T.; Myers, P. C.; Fazio, G. G.（2008 年 3 月）。“变色龙座 I 恒星形成区的盘族群”。《天体物理学杂志》，675（2）：1375--1406。arXiv:0803.1019。Bibcode:2008ApJ...675.1375L。doi:10.1086/527347。ISSN 0004-637X。  
\item Sanghi, Aniket; Liu, Michael C.; Best, William M.; Dupuy, Trent J.; Siverd, Robert J.; Zhang, Zhoujian; Hurt, Spencer A.; Magnier, Eugene A.; Aller, Kimberly M.; Deacon, Niall R.（2023 年 9 月 6 日）。“夏威夷红外视差计划 VI：利用光学至中红外 SED 并与 BT-Settl 与 ATMO~2020 模型大气比较来确定 $1000+$ 个超冷矮星与行星质量天体的基本性质”。arXiv:2309.03082 [astro-ph.SR]。  
\item Sanghi, Aniket; Liu, Michael C.; Best, William M.; Dupuy, Trent J.; Siverd, Robert J.; Zhang, Zhoujian; Hurt, Spencer A.; Magnier, Eugene A.; Aller, Kimberly M.; Deacon, Niall R.（2023 年 9 月 7 日）。“超冷天体基本性质表”。Zenodo：1。doi:10.5281/zenodo.8315643。  
\item Marsh, Kenneth A.（2010 年 2 月 1 日）。“$\rho$~蛇夫座云核心中的一颗年轻行星质量天体”。《天体物理学杂志快报》，709（2）：L158--L162。arXiv:0912.3774。Bibcode:2010ApJ...709L.158M。doi:10.1088/2041-8205/709/2/L158。S2CID 29098549。  
\item Beichman, C.; Gelino, Christopher R.; Kirkpatrick, J. Davy; Barman, Travis S.; Marsh, Kenneth A.; Cushing, Michael C.; Wright, E. L.（2013 年）。“最冷的棕矮星（或自由漂浮行星）？：Y 矮星 WISE~1828+2650”。《天体物理学杂志》，764（1）：101。arXiv:1301.1669。Bibcode:2013ApJ...764..101B。doi:10.1088/0004-637X/764/1/101。S2CID 118575478。  
\item Delorme, Philippe（2012 年 9 月 25 日）。“CFBDSIR2149$-$0403：AB~Doradus 年轻运动群中的一颗 $4$--$7$ 木星质量自由漂浮行星？”。《天文学与天体物理学》，548A：26。arXiv:1210.0305。Bibcode:2012A\&A...548A..26D。doi:10.1051/0004-6361/201219984。S2CID 50935950。  
\item Scholz, Alexander; Jayawardhana, Ray; Muzic, Koraljka; Geers, Vincent; Tamura, Motohide; Tanaka, Ichi（2012 年 9 月 1 日）。“附近年轻星团中的亚恒星天体（SONYC）VI：NGC~1333 的行星质量范围”。《天体物理学杂志》，756（1）：24。arXiv:1207.1449。Bibcode:2012ApJ...756...24S。doi:10.1088/0004-637X/756/1/24。ISSN 0004-637X。S2CID 119251742。  
\item “NAME Serpens Cluster”。simbad.cds.unistra.fr。检索于 2023 年 9 月 7 日。
\item Filippazzo, Joseph C.; Rice, Emily L.; Faherty, Jacqueline; Cruz, Kelle L.; Van Gordon, Mollie M.; Looper, Dagny L.（2015 年 9 月 1 日）。“跨越恒星到行星质量范围的年轻与场星年龄天体的基本参数与光谱能量分布”。《天体物理学杂志》，810（2）：158。arXiv:1508.01767。Bibcode:2015ApJ...810..158F。doi:10.1088/0004-637X/810/2/158。ISSN 0004-637X。S2CID 89611607。  
\item Gagné, Jonathan（2014 年 3 月 10 日）。“BANYAN II：附近年轻运动学群体中具有既往年轻性迹象的极低质量与亚恒星候选成员”。《天体物理学杂志》，783（2）：121。arXiv:1312.5864。Bibcode:2014ApJ...783..121G。doi:10.1088/0004-637X/783/2/121。S2CID 119251619。  
\item Schneider, Adam C.（2014 年 1 月 9 日）。“年轻 L 矮星 WISE~J174102.78$-$464225.5 的发现”。《天文学杂志》，147（2）：34。arXiv:1311.5941。Bibcode:2014AJ....147...34S。doi:10.1088/0004-6256/147/2/34。S2CID 38602758。  
\item Zapatero Osorio, M. R.; Lodieu, N.; Béjar, V. J. S.; Martín, Eduardo L.; Ivanov, V. D.; Bayo, A.; Boffin, H. M. J.; Muzic, K.; Minniti, D.; Beamín, J. C.（2016 年 8 月 1 日）。“WISE~J085510.74$-$071442.5 的近红外测光”。《天文学与天体物理学》，592：A80。arXiv:1605.08620。Bibcode:2016A\&A...592A..80Z。doi:10.1051/0004-6361/201628662。ISSN 0004-6361。  
\item Luhman, Kevin L.（2014 年 5 月 10 日）。“在距离太阳 2~pc 处发现一颗 $\sim250$~K 的棕矮星”。《天体物理学杂志快报》，786（2）：L18。arXiv:1404.6501。Bibcode:2014ApJ...786L..18L。doi:10.1088/2041-8205/786/2/L18。S2CID 119102654。  
\item Gagné, Jonathan; Gonzales, Eileen C.; Faherty, Jacqueline K.（2018 年）。“Gaia~DR2 的确认：2MASS~J12074836$-$3900043 为 TW~Hya 协会成员”。《美国天文学会研究札记》，2（2）：17。arXiv:1804.09625。Bibcode:2018RNAAS...2...17G。doi:10.3847/2515-5172/aac0fd。  
\item Gagné, Jonathan; Gonzales, Eileen C.; Faherty, Jacqueline K.（2018 年 5 月 1 日）。“Gaia~DR2 的确认：2MASS~J12074836--3900043 为 TW~HYA 协会成员”。《AAS 研究札记》，2（2）：17。arXiv:1804.09625。Bibcode:2018RNAAS...2...17G。doi:10.3847/2515-5172/aac0fd。ISSN 2515-5172。  
\item Gagné, Jonathan（2014 年 4 月 10 日）。“TWA 中最冷的孤立棕矮星候选成员”。《天体物理学杂志快报》，785（1）：L14。arXiv:1403.3120。Bibcode:2014ApJ...785L..14G。doi:10.1088/2041-8205/785/1/L14。S2CID 119269921。  
\item Liu, Michael C.（2016 年 12 月 9 日）。“夏威夷红外视差计划 II：年轻的超冷场矮星”。《天体物理学杂志》，833（1）：96。arXiv:1612.02426。Bibcode:2016ApJ...833...96L。doi:10.3847/1538-4357/833/1/96。S2CID 119192984。  
\item Gagné, Jonathan（2014 年 9 月 1 日）。“SIMP~J2154$-$1055：Argus 协会中的一颗新的低引力 L4$\beta$ 棕矮星候选成员”。《天体物理学杂志快报》，792（1）：L17。arXiv:1407.5344。Bibcode:2014ApJ...792L..17G。doi:10.1088/2041-8205/792/1/L17。S2CID 119118880。  
\item Kellogg, Kendra（2016 年 4 月 11 日）。“TW~Hydrae 协会中距离最近的孤立成员是类比气态巨行星的天体”。《天体物理学杂志快报》，821（1）：L15。arXiv:1603.08529。Bibcode:2016ApJ...821L..15K。doi:10.3847/2041-8205/821/1/L15。S2CID 119289711。  
\item Peña Ramírez, K.; Béjar, V. J. S.; Zapatero Osorio, M. R.（2016 年 2 月 1 日）。“上天蝎座协会中一颗新的自由漂浮行星”。《天文学与天体物理学》，586：A157。arXiv:1511.05586。Bibcode:2016A\&A...586A.157P。doi:10.1051/0004-6361/201527425。ISSN 0004-6361。S2CID 55940316。  
\item Best, William M. J.; Liu, Michael C.; Magnier, Eugene A.; Bowler, Brendan P.; Aller, Kimberly M.; Zhang, Zhoujian; Kotson, Michael C.; Burgett, W. S.; Chambers, K. C.; Draper, P. W.; Flewelling, H.; Hodapp, K. W.; Kaiser, N.; Metcalfe, N.; Wainscoat, R. J.（2017 年 3 月 1 日）。“利用 Pan-STARRS1 与 WISE 搜寻 L/T 过渡矮星 III：金牛座与天蝎—半人马座中的年轻 L 矮星发现与自行星表”。《天体物理学杂志》，837（1）：95。arXiv:1702.00789。Bibcode:2017ApJ...837...95B。doi:10.3847/1538-4357/aa5df0。ISSN 0004-637X。  
\item Zapatero Osorio, M. R.; Béjar, V. J. S.; Lodieu, N.; Manjavacas, E.（2018 年 3 月 1 日）。“确认昴宿星团中质量最低的成员”。《皇家天文学会月刊》，475（1）：139--153。arXiv:1712.01698。Bibcode:2018MNRAS.475..139Z。doi:10.1093/mnras/stx3154。ISSN 0035-8711。  
\item Gagné, Jonathan; Allers, Katelyn N.; Theissen, Christopher A.; Faherty, Jacqueline K.; Bardalez Gagliuffi, Daniella; Artigau, Étienne（2018 年 2 月 1 日）。“2MASS~J13243553+6358281 是 AB~Doradus 运动群中的早期 T 型行星质量天体”。《天体物理学杂志》，854（2）：L27。arXiv:1802.00493。Bibcode:2018ApJ...854L..27G。doi:10.3847/2041-8213/aaacfd。ISSN 0004-637X。  
\item Vos, Johanna M.; Faherty, Jacqueline K.; Gagné, Jonathan; Marley, Mark; Metchev, Stanimir; Gizis, John; Rice, Emily L.; Cruz, Kelle（2022 年）。“让伟大世界旋转：利用斯皮策空间望远镜揭示年轻巨行星类比体的风暴与湍动本性”。《天体物理学杂志》，924（2）：68。arXiv:2201.04711。Bibcode:2022ApJ...924...68V。doi:10.3847/1538-43
\item “系外行星百科全书——2MASS~J0718$-$6415”。系外行星百科全书。检索于 2021 年 1 月 31 日。  
\item Seamons, Brian; Manjavacas, Elena; Oliveros-Gómez, Natalia（2025 年 1 月）。“棕矮星还是大质量行星？：探究 WISE~J225540.75$-$311842.0 的奥秘”。美国天文学会会议摘要，245：464.09。Bibcode:2025AAS...24546409S。  
\item Schneider, Adam C.; Burgasser, Adam J.; Bruursema, Justice; Munn, Jeffrey A.; Vrba, Frederick J.; Caselden, Dan; Kabatnik, Martin; Rothermich, Austin; Sainio, Arttu; Bickle, Thomas P.; Dahm, Scott E.; Meisner, Aaron M.; Kirkpatrick, J. Davy; Suárez, Genaro; Gagné, Jonathan（2023 年 2 月 1 日）。“比红更红：一颗异常偏红的 L/T 过渡矮星的发现”。《天体物理学杂志》，943（2）：L16。arXiv:2301.02322。Bibcode:2023ApJ...943L..16S。doi:10.3847/2041-8213/acb0cd。ISSN 0004-637X。S2CID 255522681。
\item Mróz, Przemek; Poleski, Radosław; Gould, Andrew; Udalski, Andrzej; Sumi, Takahiro; Szymański, Michał K.; Soszyński, Igor; Pietrukowicz, Paweł; Kozłowski, Szymon; Skowron, Jan; Ulaczyk, Krzysztof; Albrow, Michael D.; Chung, Sun-Ju; Han, Cheongho; Hwang, Kyu-Ha; Jung, Youn Kil; Kim, Hyoun-Woo; Ryu, Yoon-Hyun; Shin, In-Gu; Shvartzvald, Yossi; Yee, Jennifer C.; Zang, Weicheng; Cha, Sang-Mok; Kim, Dong-Jin; Kim, Seung-Lee; Lee, Chung-Uk; Lee, Dong-Joo; Lee, Yongseok; Park, Byeong-Gon; 等（2020 年），“在持续时间最短的引力微透镜事件中探测到一颗类地质量的流浪行星候选体”，《天体物理学杂志》，903（1）：L11，arXiv:2009.12377，Bibcode:2020ApJ...903L..11M，doi:10.3847/2041-8213/abbfad，S2CID 221971000。  
\item Jung, Youn Kil; Hwang, Kyu-Ha; Yang, Hongjing; Gould, Andrew; Yee, Jennifer C.; Han, Cheongho; Albrow, Michael D.; Chung, Sun-Ju; Ryu, Yoon-Hyun（2024 年 5 月 27 日）。“KMT-2023-BLG-2669：第九个具有 $\theta_{\rm E}$ 测量的自由漂浮行星候选体”。arXiv:2405.16857 [astro-ph.EP]。  
\item Mróz, Przemek; Udalski, Andrzej; Bennett, David P.; Ryu, Yoon-Hyun; Sumi, Takahiro; Shvartzvald, Yossi; Skowron, Jan; Poleski, Radosław; Pietrukowicz, Paweł; Kozłowski, Szymon; Szymański, Michał K.; Wyrzykowski, Łukasz; Soszyński, Igor; Ulaczyk, Krzysztof; Rybicki, Krzysztof（2019 年 2 月 1 日）。“通过微透镜发现两颗新的自由漂浮或宽轨道行星”。《天文学与天体物理学》，622：A201。arXiv:1811.00441。Bibcode:2019A\&A...622A.201M。doi:10.1051/0004-6361/201834557。ISSN 0004-6361。  
\item Becky Ferreira（2018 年 11 月 9 日）。“罕见目击：两颗不绕恒星运行的流浪行星”。Motherboard。检索于 2019 年 2 月 10 日。  
\item Jake Parks（2018 年 11 月 16 日）。“这两颗新的‘流浪行星’在没有恒星的情况下于宇宙中游荡”。《Discover》杂志。原文存档于 2018 年 11 月 16 日。检索于 2019 年 2 月 10 日。  
\item Jake Parks（2018 年 11 月 15 日）。“发现两颗自由放养的行星在银河系中孤独漫游”。《Astronomy》杂志。检索于 2019 年 2 月 10 日。  
\item “系外行星目录”。Exoplanet Exploration: Planets Beyond our Solar System。2018 年 11 月 10 日。检索于 2021 年 1 月 4 日。  
\item Miyazaki, S.; Sumi, T.; Bennett, D. P.; Gould, A.; Udalski, A.; Bond, I. A.; Koshimoto, N.; Nagakane, M.; Rattenbury, N.; Abe, F.; Bhattacharya, A.; Barry, R.; Donachie, M.; Fukui, A.; Hirao, Y.（2018 年 9 月 1 日）。“MOA-2015-BLG-337：一个宿主位于低质量棕矮星/行星边界的行星系统，或一对棕矮星双星”。《天文学杂志》，156（3）：136。arXiv:1804.00830。Bibcode:2018AJ....156..136M。doi:10.3847/1538-3881/aad5ee。ISSN 0004-6256。  
\item Kim, Hyoun-Woo; Hwang, Kyu-Ha; Gould, Andrew; Yee, Jennifer C.; Ryu, Yoon-Hyun; Albrow, Michael D.; Chung, Sun-Ju; Han, Cheongho; Jung, Youn Kil; Lee, Chung-Uk; Shin, In-Gu; Shvartzvald, Yossi; Zang, Weicheng; Cha, Sang-Mok; Kim, Dong-Jin; Kim, Seung-Lee; Lee, Dong-Joo; Lee, Yongseok; Park, Byeong-Gon; Pogge, Richard W.（2021 年）。“KMT-2019-BLG-2073：第四个满足 $\theta_{\rm E}<10\,\mu\mathrm{as}$ 的自由漂浮行星候选体”。《天文学杂志》，162（1）：15。arXiv:2007.06870。Bibcode:2021AJ....162...15K。doi:10.3847/1538-3881/abfc4a。S2CID 235445277。  
\item Mróz, Przemek; 等（2020 年），“微透镜事件 OGLE-2019-BLG-0551 中的一颗自由漂浮或宽轨道行星”，《天文学杂志》，159（6）：262，arXiv:2003.01126，Bibcode:2020AJ....159..262M，doi:10.3847/1538-3881/ab8aeb，S2CID 211817861。  
\item Kaczmarek, Zofia; McGill, Peter; Evans, N. Wyn; Smith, Leigh C.; Wyrzykowski, Łukasz; Howil, Kornel; Jabłońska, Maja（2022 年 8 月 1 日）。“尘埃中的暗透镜：VVV 的视差微透镜事件”。《皇家天文学会月刊》，514（4）：4845--4860。arXiv:2205.07922。Bibcode:2022MNRAS.514.4845K。doi:10.1093/mnras/stac1507。ISSN 0035-8711。 
\item Koshimoto, Naoki; Sumi, Takahiro; Bennett, David P.; Bozza, Valerio; Mróz, Przemek; Udalski, Andrzej; Rattenbury, Nicholas J.; Abe, Fumio; Barry, Richard; Bhattacharya, Aparna; Bond, Ian A.; Fujii, Hirosane; Fukui, Akihiko; Hamada, Ryusei; Hirao, Yuki（2023 年 3 月 14 日）。“基于 MOA-II 为期 9 年的银河系核球巡天得到的类地质量与海王星质量自由漂浮行星候选体”。《天文学杂志》，166（3）：107。arXiv:2303.08279。Bibcode:2023AJ....166..107K。doi:10.3847/1538-3881/ace689。  
\item Shah, Sarang（2019 年）。“利用最新一代望远镜对微透镜事件的分析”。坎特伯雷大学——经由 UC Research Repository。  
\item Kenworthy, M. A.; Klaassen, P. D.; 等（2020 年 1 月）。“指向年轻外环凌星系统 J1407（V1400~Cen）的 ALMA 与 NACO 观测”。《天文学与天体物理学》，633：A115。arXiv:1912.03314。Bibcode:2020A\&A...633A.115K。doi:10.1051/0004-6361/201936141。  
\item “流浪行星”。warwick.ac.uk。检索于 2025 年 4 月 26 日。  
\item Nicoll, James Davis（2020 年 4 月 30 日）。“远离任何恒星：关于流浪世界的五个故事”。Reactor。检索于 2025 年 4 月 26 日。  
\item Langford, David; Stableford, Brian（2021 年）。“外行星”。载于 Clute, John; Langford, David; Sleight, Graham（编），《科幻百科全书》（第 4 版）。检索于 2023 年 6 月 13 日。数十年来，冥王星因其作为太阳系表观的 Ultima Thule 而受到一定程度的特别关注……冥王星在其轨道似乎标记了太阳系最外边界的时期，正是因此而广受欢迎。  
\item Womack, Lacey（2020 年 5 月 3 日）。“YouTube 上 15 个最好的视频式 ARG”。ScreenRant。检索于 2025 年 4 月 26 日。
\end{enumerate}
\subsection{参考书目}
\begin{itemize}
\item Stevenson 撰写的文章《星际空间中可维持生命行星的可能性》，内容与《Nature》上的相关文章相似，但信息更加详尽。  
\end{itemize}
\subsection{外部链接}
\begin{itemize}
\item “行星”的定义（IAU 第 B5 号决议）  
\item 《奇异的新世界或可形成微型太阳系》，Robert Roy Britt，SPACE.com，2006 年 6 月 5 日 上午 11:35（ET）  
\item “行星”与“冥王类天体（plutons）”的 IAU 草案定义新闻稿（国际天文学联合会，2006 年）
\end{itemize}

